% === Begin Label Data ===
% ID: 167
% 难度星级: 
% 标签: 
% 备注: 
% 组卷引用次数: 0
% === End  Label Data ===

\begin{problem}{2019}{G}{课标I卷（理）}{21}{概率}
为治疗某种疾病，研制了甲、乙两种新药，希望知道哪种新药更有效，为此进行动物实验. 试验方案如下：每一轮选取两只白鼠对药效进行对比试验. 对于两只白鼠，随机选一只施以甲药，另一只施以乙药. 一轮的疗效结果得出后，再安排下一轮试验. 当其中一种药治愈的白鼠比另一种药治愈的白鼠多 4 只时，就停止试验，并认为治愈只数多的药更有效. 为了方便描述问题，约定：对于每轮试验，若施以甲药的白鼠治愈且施以乙药的白鼠未治愈则甲药得 1 分，乙药得 -1 分；若施以乙药的白鼠治愈且施以甲药的白鼠未治愈则乙药得 1 分，甲药得 -1 分；若都治愈或都未治愈则两种药均得 0 分. 甲、乙两种药的治愈率分别记为 $\alpha$ 和 $\beta$，一轮试验中甲药的得分记为 $X$.

（1）求 $X$ 的分布列;

（2）若甲药、乙药在试验开始时都赋予 4 分， $p_i(i=0,1,\cdots,8)$ 表示“甲药的累计得分为 $i$ 时，最终认为甲药比乙药更有效”的概率，则 $p_0=0, p_8=1, p_i = ap_{i-1} + b p_i + c p_{i+1}(i=1,2,\cdots,7)$，其中 $a=P(X=-1), b=P(X=0), c=P(X=1)$. 假设 $\alpha=0.5, \beta=0.8$.

(i) 证明：$\{p_{i+1}-p_i\}(i=0,1,2,\cdots,7)$ 为等比数列;

(ii) 求 $p_4$，并根据 $p_4$ 的值解释这种试验方案的合理性.

\end{problem}


\begin{solutions}
(1) $X$ 的所有可能取值为 $-1, 0, 1$. 

$P(X=-1) = (1-\alpha)\beta, P(X=0) = \alpha\beta + (1-\alpha)(1-\beta), P(X=1) = \alpha(1-\beta)$.

所以 $X$ 的分布列为
\begin{center}
\begin{tabular}{|c|c|c|c|}
\hline
$X$ & $-1$ & $0$ & $1$ \\
\hline
$P$ & $(1-\alpha)\beta$ & $\alpha\beta + (1-\alpha)(1-\beta)$ & $\alpha(1-\beta)$ \\
\hline
\end{tabular}
\end{center}

(2) \circled{1} 由 (1) 得 $a=0.4, b=0.5, c=0.1$. 

因此 $p_i = 0.4p_{i-1} + 0.5p_i + 0.1p_{i+1}$, 故 $0.1(p_{i+1}-p_i) = 0.4(p_i-p_{i-1})$, 即 $p_{i+1}-p_i = 4(p_i-p_{i-1})$. 

又因为 $p_1-p_0 = p_1 \ne 0$, 所以 $\{p_{i+1}-p_i\} (i=0, 1, 2, \cdots, 7)$ 为公比为 4, 首项为 $p_1$ 的等比数列.

\circled{2} 由 \circled{1} 可得 $p_8 = p_8-p_7+p_7-p_6+\cdots+p_1-p_0+p_0 = (p_8-p_7)+(p_7-p_6)+\cdots+(p_1-p_0) = \dfrac{4^8-1}{3}p_1$.

由于 $p_8=1$, 故 $p_1=\dfrac{3}{4^8-1}$, 所以 $p_4 = (p_4-p_3)+(p_3-p_2)+(p_2-p_1)+(p_1-p_0) = \dfrac{4^4-1}{3}p_1 = \dfrac{1}{257}$.

$p_4$ 表示最终认为甲药更有效的概率. 由计算结果可以看出, 在甲药治愈率为 $0.5$, 乙药治愈率为 $0.8$ 时, 认为甲药更有效的概率为 $p_4 = \dfrac{1}{257} \approx 0.0039$, 此时得出错误结论的概率非常小, 说明这种试验方案合理.


\end{solutions}